% The frames named in sec:reference_frames, in one schematic elevation.
% The arm is drawn only to place the frames; it is not a kinematic diagram and
% carries no link lengths. The desired frame is drawn offset from the
% end-effector frame so that the position error used by the impedance law is a
% visible displacement rather than only an equation.
%
% Uses only tikz + arrows.meta/calc, which config/packages.tex loads. Colour
% follows the established meaning: grey for measured or structural quantities,
% orange for the commanded frame that a parameter selects, blue for the surface
% frame the task is defined against.
%
% Clearance note. The end-effector cluster is the crowded part of this drawing:
% four labels and the error arrow meet there. Three directions are kept free
% around it and nothing else may be placed on them: up-right for the error
% arrow to {d}, down-right for z_EE toward the surface, and up-left for the
% frame label itself.
\begin{tikzpicture}[
    font=\small,
    link/.style={draw=gray!55, line width=2.4pt, line cap=round},
    base/.style={draw=gray!70!black, fill=gray!22},
    surf/.style={draw=gray!70!black, fill=gray!12},
    meas/.style={-{Latex[length=1.8mm]}, gray!55!black},
    des/.style={-{Latex[length=1.8mm]}, densely dashed, orange!75!black},
    tsk/.style={-{Latex[length=1.8mm]}, blue!55!black},
    err/.style={-{Latex[length=2.0mm]}, thick, gray!45!black},
    lbl/.style={font=\footnotesize, inner sep=1.8pt},
    jt/.style={circle, draw=gray!60!black, fill=white, inner sep=0pt,
               minimum size=3.4pt},
  ]

% ============================== base frame ===============================
\draw[base] (-0.55,-0.30) rectangle (0.55,0.10);
\node[lbl, gray!55!black] at (0,-0.54) {$\{0\}$};

% The into-page axis marks the frame origin, so the two drawn axes start
% clear of its circle rather than from the point beneath it.
\draw[meas] (0.19,0.10) -- (1.05,0.10) node[lbl, right] {$x_0$};
\draw[meas] (0,0.29) -- (0,1.20) node[lbl, above] {$z_0$};

% ============================ schematic arm ==============================
% Joint frames are indexed along this chain; the controller takes its poses
% from the robot model rather than from a drawn parameterisation.
\draw[link] (0,0.10) -- (0.95,2.55);
\draw[link] (0.95,2.55) -- (3.30,3.15);
\draw[link] (3.30,3.15) -- (4.45,1.85);
\node[jt] at (0.95,2.55) {};
\node[jt] at (3.30,3.15) {};

% The base origin marker is drawn after the chain so the first link passes
% behind it rather than through the symbol.
\draw[gray!60!black, fill=white] (0,0.10) circle (0.13);
\draw[gray!60!black] (-0.09,0.19) -- (0.09,0.01);
\draw[gray!60!black] (-0.09,0.01) -- (0.09,0.19);
\node[lbl, anchor=south east, gray!60!black] at (-0.14,0.16) {$y_0$};
\node[lbl, gray!55!black] at (1.80,3.12) {$\{i\}$};

% ========================= end-effector frame ============================
\coordinate (EE) at (4.45,1.85);
\node[jt] at (EE) {};
\draw[meas] (EE) -- ++(0.86,0.34);
\node[lbl, anchor=west, gray!55!black] at ($(EE)+(0.94,0.20)$)
      {$x_{\mathrm{EE}}$};
\draw[meas] (EE) -- ++(0.34,-0.86);
\node[lbl, anchor=east, gray!55!black] at ($(EE)+(0.10,-0.78)$)
      {$z_{\mathrm{EE}}$};
\node[lbl, anchor=east, gray!55!black] at ($(EE)+(-0.14,-0.04)$)
      {$\{\mathrm{EE}\}$};

% ========================== desired frame {d} ============================
\coordinate (D) at (5.90,2.95);
\node[jt] at (D) {};
\draw[des] (D) -- ++(0.86,0.34);
\draw[des] (D) -- ++(0.34,-0.86);
\node[lbl, anchor=south east, orange!65!black] at ($(D)+(-0.06,0.04)$)
      {$\{d\}$};

% ---- the position error the impedance law acts on -----------------------
\draw[err] (EE) -- (D);
\node[lbl, anchor=south east, gray!45!black]
      at ($(EE)!0.55!(D)+(0.02,0.04)$) {$e_p$};

% ============================ surface frame ==============================
\begin{scope}[shift={(5.95,-0.10)}, rotate=-7]
  \fill[surf] (-1.95,-0.30) rectangle (1.95,0);
  \draw[gray!70!black] (-1.95,0) -- (1.95,0);

  % t_2 points into the page and marks p_s, so the two drawn tangent-frame
  % axes start clear of its circle.
  \coordinate (S) at (-0.55,0);
  \draw[tsk] ($(S)+(0.19,0)$) -- ++(0.96,0) node[lbl, right] {$t_1$};
  \draw[tsk] ($(S)+(0,0.19)$) -- ++(0,0.86) node[lbl, above] {$n_s$};
  \draw[blue!55!black, fill=white] (S) circle (0.13);
  \draw[blue!55!black] ($(S)+(-0.09,0.09)$) -- ($(S)+(0.09,-0.09)$);
  \draw[blue!55!black] ($(S)+(-0.09,-0.09)$) -- ($(S)+(0.09,0.09)$);
  \node[lbl, anchor=south east, blue!50!black] at ($(S)+(-0.14,0.16)$) {$t_2$};
  \node[lbl, anchor=north, blue!50!black] at ($(S)+(0,-0.42)$) {$p_s$};
  \node[lbl, blue!50!black] at (1.55,0.62) {$\{S\}$};
\end{scope}

\end{tikzpicture}
